% !TeX root = project-report-basic-cn.tex
% !TeX program = lualatex
% Adapted from Downloads/project-report-cn-v4.pdf and its matching source at bb1d4cd.
% Historical result values are retained; no new event access or scientific execution.
\documentclass[11pt,a4paper,fontset=fandol,scheme=plain]{ctexart}
\usepackage[margin=21mm,headheight=14pt]{geometry}
\usepackage{lmodern,amsmath,amssymb,graphicx,booktabs,array,tabularx,tikz}
\usepackage{xcolor,enumitem,fancyhdr,caption,xurl,float}
\usepackage[colorlinks=true,linkcolor=teal,citecolor=teal,urlcolor=teal]{hyperref}
\definecolor{reportblue}{HTML}{18536B}
\hypersetup{pdftitle={从输入到质量谱：神经网络筛选的五个步骤},pdfauthor={HiggsML 项目}}
\pagestyle{fancy}
\fancyhf{}
\fancyhead[L]{\small\textcolor{reportblue}{HiggsML | 神经网络筛选基础报告}}
\fancyhead[R]{\small 2026 年 11 月 4 日}
\fancyfoot[C]{\thepage}
\setlength{\parindent}{0pt}
\setlength{\parskip}{7pt}
\setlength{\emergencystretch}{2em}
\setlist[itemize]{leftmargin=1.5em,itemsep=3pt,topsep=4pt}
\setlist[enumerate]{leftmargin=1.7em,itemsep=4pt,topsep=4pt}
\captionsetup{font=small,labelfont=bf,justification=raggedright,singlelinecheck=false}
\captionsetup[table]{name=表}
\captionsetup[figure]{name=图}
\newcolumntype{Y}{>{\raggedright\arraybackslash}X}
\newcommand{\reportsection}[1]{\section*{\textcolor{reportblue}{#1}}}
\newcommand{\gev}{\,\mathrm{GeV}}
\newcommand{\takeaway}[1]{\par\medskip\noindent\colorbox{reportblue!7}{\parbox{\dimexpr\linewidth-2\fboxsep\relax}{\textbf{本步得到什么：}#1}}\par}
% Display the unchanged upper ABCD panels only. Restore numerical x ticks at
% the positions read from the original lower panels; no curve is regenerated.
\newcommand{\abcdplot}[3]{%
\resizebox{\linewidth}{!}{%
\begin{tikzpicture}[x=1bp,y=1bp]
\node[anchor=south west,inner sep=0] at (0,0)
 {\includegraphics[width=792bp,viewport=0 295 792 576,clip]{#1}};
\foreach \x/\lab in {#2}{\node[anchor=north,font=\fontsize{10}{12}\selectfont,inner sep=0] at (\x,-3) {\lab};}
\foreach \x/\lab in {#3}{\node[anchor=north,font=\fontsize{10}{12}\selectfont,inner sep=0] at (\x,-3) {\lab};}
\node[anchor=north,font=\fontsize{11}{13}\selectfont] at (220,-20) {Four-lepton mass [GeV]};
\node[anchor=north,font=\fontsize{11}{13}\selectfont] at (610,-20) {Four-lepton mass [GeV]};
\end{tikzpicture}}}

\begin{document}
{\LARGE\bfseries\color{reportblue} 从输入到质量谱\par}
\vspace{4pt}
{\Large 一个神经网络筛选的五个步骤\par}
\vspace{4pt}
本报告以下载文件 \texttt{project-report-cn-v4.pdf} 为基础，只讲清楚一条流程：
\begin{center}
\textbf{输入是什么 $\to$ 怎样学出分数 $\to$ 怎样定阈值}\par
\textbf{$\to$ 留下哪些事件 $\to$ 质量谱怎样变化。}
\end{center}
全篇使用原报告中的 \textbf{ABCD、随机种子 42}。这里的数值与图沿用既有 draft001 结果记录，
本次仅改写解释与展示；结果来源及本机运行的区别见末页。

\reportsection{1. 输入是什么？}
\textbf{先认识一条事件。} 一条 MC 事件记录包含重建轻子的动量、能量、电荷和类型等信息。
信号来自模拟的 $H\to ZZ^*\to4\ell$；背景来自另一份四轻子 MC 样本。
本项目只分析两个电子和两个缪子（$2e2\mu$），并要求 $105\leq m_{4\ell}<140\gev$。
样本来自 ATLAS 2020 教育用途 13 TeV 集合\cite{OpenData}，没有使用真实碰撞数据。

\textbf{先定义共同基线，再让网络打分。} 事件首先通过触发、轻子品质、四轻子配对及动量和质量要求。
后文的“筛选前”已经通过这些条件；“筛选后”再加一条网络高分要求。
共同筛选的完整条件和来源保留在附录 A。

\textbf{把每条事件整理成 19 个数字。} 四个轻子按横动量 $p_T$ 从高到低编号。
电子对和缪子对中，质量更接近 $91.1876\gev$ 的一对记作 $Z_1$，另一对记作 $Z_2$。
ABCD 表示同时使用以下四组量：
\begin{table}[H]
\caption{同一事件的四组输入。各组共同构成一个 19 维向量。}
\small
\renewcommand{\arraystretch}{1.18}
\begin{tabularx}{\linewidth}{@{}p{.06\linewidth}p{.38\linewidth}cY@{}}
\toprule
组 & 输入 & 个数 & 直观含义 \\
\midrule
A & 四个轻子的 $p_T$ 和赝快度 $\eta$ & 8 & 各轻子的动量尺度与相对束流的方向。\\
B & $m_{Z1},m_{Z2},\Delta R_{Z1},\Delta R_{Z2}$ & 4 & 两个轻子对的质量和各对内部的角分离。\\
C & $p_{T,4\ell},|\Delta\phi_{ZZ}|$ & 2 & 四轻子系统及两个轻子对的整体运动学。\\
D & $\cos\theta^*,\cos\theta_1,\cos\theta_2$；$\phi_{\rm decay},\phi_{\rm prod}$ & 5 & 三个方向角余弦和两个平面角，描述产生与衰变的几何关系。\\
\bottomrule
\end{tabularx}
\end{table}
这些量由公开的轻子字段取得或计算。$m_{4\ell}$ 由四个轻子的总四动量计算，
用于共同质量窗口和最后的谱图，\textbf{不作为分类器的显式输入}。
但其他输入与质量相关，不能据此认为网络选择与质量无关。
\takeaway{每个基线事件有一个 19 维输入向量；后面训练的是从这 19 个数字到一个分数的映射。}

\clearpage
\reportsection{2. 网络怎样学出分数？}
\textbf{先分清三份样本的职责。} 共同基线按物理事件组划分，同组事件不会跨子集。
60/20/20 是分配概率，实际数量不必恰好等于这些比例。
\begin{table}[H]
\caption{原报告记录的三个子集。信号、背景列为 MC 物理事件组数。}
\small
\renewcommand{\arraystretch}{1.2}
\begin{tabularx}{\linewidth}{@{}lrrrY@{}}
\toprule
子集 & 概率 & 信号 & 背景 & 用途 \\
\midrule
训练 & 60\% & 27,026 & 2,175 & 学习标准化参数和网络参数。\\
验证 & 20\% & 9,008 & 720 & 选择检查点，并确定分数阈值。\\
最终评估 & 20\% & 9,049 & 698 & 应用冻结的模型与阈值，计算结果。\\
\bottomrule
\end{tabularx}
\end{table}

\textbf{第一步：统一输入尺度。} 只用训练集计算每一列的均值和标准差，将该列变为
\[
 \widetilde{x}_j=\frac{x_j-\overline{x}_{j,\rm train}}{\sigma_{j,\rm train}}.
\]
这里 $j$ 表示第几个输入。验证和评估沿用相同变换，不重新估计均值和标准差。

\textbf{第二步：从输入得到预测。} 信号标签设为 $y=1$，背景为 $y=0$。
网络依次组合输入，结构为
\[
 19\longrightarrow64\longrightarrow64\longrightarrow32\longrightarrow1.
\]
中间三个数字是三个隐藏层的单元数。隐藏层使用 LayerNorm 稳定层内数值尺度、SiLU 引入非线性；
前两个隐藏层后的 Dropout 在训练时以 0.1 的概率随机屏蔽输出，以减少对固定连接的依赖。
最后用 sigmoid 将输出变成 $0$ 到 $1$ 的分数 $s$。
高分表示更倾向于信号，\textbf{分数不是已校准的物理信号概率}。

\textbf{第三步：根据错误调整参数。} 每一批训练事件都计算加权二元交叉熵：
\[
 \mathcal L_{\rm train}\ \propto\ -\sum_i a_i\left[y_i\log s_i+(1-y_i)\log(1-s_i)\right].
\]
信号得到很低分、背景得到很高分，都会增加损失。反向传播计算参数应怎样变化，AdamW 据此更新参数。
$a_i$ 是非负训练权重：取物理权重的绝对值，再除以所属训练类别的平均绝对权重。
训练权重用于学习；第 4 节的物理产额另用带符号权重。

\textbf{第四步：用验证集选定检查点。} 每轮计算验证 AUC，检查信号是否通常排在背景之前。
0.5 对应随机排序，1 对应完全区分。AUC 比此前最佳值提高超过 $10^{-4}$ 才保存新检查点；
连续 20 轮未更新就停止，最多 200 轮。
原报告中 ABCD、seed 42 在第 87 轮停止，最终使用\textbf{第 67 轮}的检查点。
训练批量为 1024，学习率 $10^{-3}$，权重衰减 $10^{-4}$，使用 CPU float32。
\takeaway{现在已经得到一个固定网络：给它一条事件的输入，就能算出分数；还需要下一步的阈值，才能决定是否保留。}

\clearpage
\reportsection{3. 阈值怎样确定？}
\textbf{分数需要配合一条选择规则。} 本报告把满足 $s\geq t$ 的事件称为高分事件，
也就是谱图中的“网络筛选后”；$s<t$ 的事件属于低分类别。$t$ 就是阈值。

\textbf{阈值只由验证集背景确定。} 使用上一页选好的网络，为验证集背景打分，按下面的固定规则计算：
\begin{enumerate}
\item 将背景事件按分数从小到大排序。
\item 累加每条事件的\textbf{绝对物理权重}。
\item 第一次达到总绝对权重一半时，取该位置事件的实际分数作为 $t$。
\end{enumerate}
这就是加权中位数规则。它没有用最终评估谱图来寻找看起来最好的筛选。

\textbf{一个算得完的例子。} 下表是\textbf{教学示例，不是实际运行数据}：
\begin{table}[H]
\caption{排序、累加绝对权重，再确定阈值。}
\centering
\renewcommand{\arraystretch}{1.3}
\begin{tabular}{lrrrrr}
\toprule
排序后的分数 & 0.10 & 0.30 & 0.60 & 0.80 & 0.90\\
绝对权重 & 1 & 2 & 1 & 4 & 2\\
累计绝对权重 & 1 & 3 & 4 & \textbf{8} & 10\\
\bottomrule
\end{tabular}
\end{table}
总权重为 10，一半是 5；累计到 0.60 时只有 4，累计到 0.80 时达到 8，第一次越过 5。
因此 $t=0.80$，分数 0.80、0.90 进入高分类别。
它们的绝对权重占比是 $(4+2)/10=60\%$。离散权重和同分事件使结果不必恰好为 50\%。

\textbf{回到原报告的实际模型。} ABCD、seed 42 的阈值是
\[
 \boxed{t\simeq0.746724},\qquad
 \text{高分事件：}\ s\geq t.
\]
这里展示的是舍入值，既有评估使用完整精度。所有等于阈值的事件都进入高分类别，不拆分同分事件。
验证背景的加权中位数规则，也不保证另一份评估样本的带符号背景产额恰好保留一半。

\textbf{冻结之后再评估。} 顺序是：保存输入标准化参数、固定网络检查点、固定阈值，
然后才对最终评估事件应用这条选择。评估结果不再用于调整模型或阈值。
所谓“冻结”，就是不根据评估图形是否好看再改这些选择。

\takeaway{本报告的“筛选后”有明确含义：用已冻结的 ABCD 网络打分，再保留达到或超过既定阈值的事件。}

\clearpage
\reportsection{4. 留下哪些事件？}
先对同一最终评估样本打分，再比较基线与其中的高分子集。
\begin{table}[H]
\caption{ABCD、seed 42 的 MC 物理事件组数，沿用原报告。}
\centering
\renewcommand{\arraystretch}{1.25}
\begin{tabular}{lrr}
\toprule
过程 & 筛选前 & 筛选后（高分）\\
\midrule
信号 & 9,049 & 8,761\\
背景 & 698 & 386\\
\bottomrule
\end{tabular}
\end{table}
高分事件来自原来的基线事件，不是一批新生成或独立的样本。

\textbf{为什么九千多个模拟事件只对应约 2.77 个预期事件？}
MC 事件数表示模拟采样的规模；实际预期产额通过物理权重换算，二者不能直接等同。
本项目按 $10\,\mathrm{fb}^{-1}$ 的积分亮度归一化，单条事件的物理权重为
\[
 w_i^{\rm phys}=\frac{\mathcal L\,\sigma_{\rm eff}}{S_{\rm gen}}w_i^{\rm gen}.
\]
$\mathcal L$ 是目标积分亮度，$\sigma_{\rm eff}$ 是样本的有效截面，
$S_{\rm gen}$ 是归一化用的生成器权重和，$w_i^{\rm gen}$ 是该事件的生成器权重。
它可以带负号；负贡献保留原符号。样本常数与适用性说明见附录 B。

\textbf{为什么还要除以 0.2？} 最终评估只接收约 20\% 的共同基线事件。
为估计相同目标亮度下的产额，评估时使用 $u_i=w_i^{\rm phys}/0.2$，再求和：
\[
 Y_{\rm before}=\sum_{i\in E}u_i,\qquad
 Y_{\rm after}=\sum_{i\in E,\,s_i\geq t}u_i,\qquad
 R=\frac{Y_{\rm after}}{Y_{\rm before}}.
\]
$E$ 是某一过程的评估基线，$R$ 是该过程的产额保留率。
这个修正对应评估抽样概率，不能补回被物理条件排除的事件，也不增加亮度。

\begin{table}[H]
\caption{$10\,\mathrm{fb}^{-1}$ 下的带符号预期产额。误差沿用原报告的固定网络、固定阈值 MC 标准误差；保留率误差包含筛选前后的相关性。}
\centering
\renewcommand{\arraystretch}{1.25}
\begin{tabular}{lrrr}
\toprule
过程 & 筛选前产额 & 筛选后产额 & 保留率（\%）\\
\midrule
信号 & $2.7710\pm0.0293$ & $2.6824\pm0.0288$ & $96.80\pm0.19$\\
背景 & $2.3553\pm0.1683$ & $1.2027\pm0.1202$ & $51.06\pm3.57$\\
\bottomrule
\end{tabular}
\end{table}
例如 $2.6824/2.7710\simeq96.8\%$；这是产额之比，区别于按事件组计算的 $8761/9049$。
表中用原始精度计算的保留率沿用既有记录，显示值已舍入。
\takeaway{在这一固定选择下，保留约 97\% 的信号产额、约 51\% 的背景产额，即损失约 3\% 信号，同时去除约 49\% 背景。}

\clearpage
\reportsection{5. 质量谱怎样变化？先看绝对产额}
\textbf{对保留下来的事件，画它们原有的质量。} 将四个轻子的四动量相加得到 $Q$，
按 $m_{4\ell}=\sqrt{Q^2}$ 计算质量。网络只决定哪些事件进入高分子集，没有改写单个事件的质量。
在 $105$ 至 $140\gev$ 内，每 $5\gev$ 一个箱，对箱内事件的评估权重 $u_i$ 求和。

\begin{figure}[H]
\centering
\abcdplot{figures/draft001/mass-yields.pdf}
{64.909/105,109.221/110,153.532/115,197.843/120,242.155/125,286.466/130,330.778/135,375.089/140}
{455.509/105,499.821/110,544.132/115,588.443/120,632.755/125,677.066/130,721.378/135,765.689/140}
\caption{原报告图 1 的 ABCD、seed 42 上排面板。左为信号，右为背景；蓝色为共同基线，橙色为高分子集。
每点代表一个 $5\gev$ 箱，纵轴是 $10\,\mathrm{fb}^{-1}$ 下的带符号预期产额，左右纵轴刻度不同。
仅截取原图面板并补上相同的横轴刻度数字，曲线、误差棒与阴影均保留。连线只作视觉引导。}
\end{figure}

\textbf{先看信号。} 信号主要集中在 $125\gev$ 附近。筛选前后曲线相近，符合上一页“保留大部分信号”的产额结果。
这个峰是所用模拟样本及重建、选择共同形成的分布，不能由粗分箱判断质量分辨率得到改善。

\textbf{再看背景。} 背景在筛选前就分布不均匀。筛选后的总产额下降，但各质量区域压低的程度不同，
高质量尾部尤其明显。总背景减少，并不意味着原来的谱形只是整体乘了一个常数。

\textbf{误差棒和阴影怎样读？} 误差棒反映固定网络、固定阈值下的有限 MC 统计误差。
阴影表示至少一条曲线在该箱统计支持不足，属于诊断标记；若干背景箱和信号高质量尾部存在这种情况。
ABCD 高分背景的 $[105,110)\gev$ 箱为负，这是带符号权重求和的结果，不表示出现了负的实际事件数。

原报告的阴影规则为：无事件组；带符号产额 $Y$ 或分组方差 $V$ 非正；
有效计数 $Y^2/V<20$；或抵消比 $|Y|/\sum_i|u_i|<0.2$。
这里 $V$ 是各物理事件组产额贡献的平方和。该标记不是置信带，阴影箱和负值均未被删去。

\takeaway{这一页同时看到信号大体保留、背景减少；下一页去掉总量差异，单独检查形状。}

\clearpage
\reportsection{5. 质量谱怎样变化？再看归一化形状}
把每条曲线除以它自己的正的带符号总产额：若第 $b$ 个质量箱的产额为 $Y_b$，
图中显示 $f_b=Y_b/\sum_bY_b$。这样比较的是各质量区域所占的相对份额。

\begin{figure}[H]
\centering
\abcdplot{figures/draft001/mass-shapes.pdf}
{63.388/105,107.465/110,151.543/115,195.62/120,239.698/125,283.776/130,327.853/135,371.931/140}
{457.228/105,501.305/110,545.383/115,589.46/120,633.538/125,677.616/130,721.693/135,765.771/140}
\caption{原报告图 2 的 ABCD、seed 42 上排面板，采用与上一页相同的面板截取方式。
每条曲线分别归一化；误差包含共享归一化的相关性，筛选前后也共享事件。
带符号归一化后的箱值可以为负，不应直接当成概率密度。}
\end{figure}

\begin{table}[H]
\caption{原报告中的固定质量区域产额占比（\%）。这些是描述性点估计，未列不确定度，不能据此判定统计显著性。}
\centering
\renewcommand{\arraystretch}{1.2}
\begin{tabular}{llrr}
\toprule
过程 & 质量区域（GeV） & 筛选前 & ABCD 筛选后\\
\midrule
信号 & $[120,130)$ & 87.9 & 88.6\\
背景 & $[120,130)$ & 32.5 & 51.8\\
背景 & $[135,140)$ & 22.0 & 6.5\\
\bottomrule
\end{tabular}
\end{table}

\textbf{背景占比上升，为什么总背景还在减少？} 占比的分母也改变了。
用已舍入的数值作近似说明：$[120,130)\gev$ 的背景产额筛选前约为 $2.3553\times32.5\%\simeq0.77$，
筛选后约为 $1.2027\times51.8\%\simeq0.62$。
该区域的绝对背景仍下降，只是其他区域下降得更多，所以它在剩余背景中的占比上升。

\textbf{这条流程能够支持什么？} 对这份 MC 样本，冻结的 ABCD 选择保留大部分信号，压低约一半背景，
并改变背景质量谱形。因此，“没有显式输入质量”不等于“筛选与质量无关”。
现有结果属于探索性、非独立的教育与技术研究；图中误差不包含重新训练或阈值估计的全部不确定性。
这份报告不据此判断输入组合优劣、发现新共振或确认物理测量精度提升。
\takeaway{现在可以从 19 个输入一路追到最终谱图：模型给分，固定阈值选择事件，带符号权重决定产额，事件子集决定谱形。}

\clearpage
\reportsection{附录 A：共同选择及其来源}
下面保留下载版的共同基线，便于核查“筛选前”究竟包括什么。动量、质量单位均为 GeV。
\begin{table}[H]
\small
\renewcommand{\arraystretch}{1.18}
\begin{tabularx}{\linewidth}{@{}p{.22\linewidth}Y@{}}
\toprule
步骤 & 条件及目的\\
\midrule
事件触发 & \texttt{trigE} 或 \texttt{trigM} 为真。\\
轻子味道与鉴别 & $|\texttt{lep\_type}|\in\{11,13\}$，\texttt{lep\_isTightID=true}。以下品质条件依次执行，每步至少保留四个轻子。\\
隔离 & $p_T>0$；依次要求 \texttt{lep\_ptcone30}/\texttt{lep\_pt}$<0.3$ 和 \texttt{lep\_etcone20}/\texttt{lep\_pt}$<0.3$。\\
碰撞参数 & $|d_0/\sigma(d_0)|<5$（电子）或 3（缪子）；$|z_0/\cosh\eta|<0.5$ mm。\\
多重数与匹配 & 品质筛选后恰好四个轻子，至少一个 \texttt{lep\_trigMatched=true}；原始轻子数可以大于四。\\
动量与接受度 & 按 $p_T$ 降序依次 $\geq20,15,10,7$；电子 $|\eta|<2.47$、缪子 $|\eta|<2.7$。\\
电荷与配对 & 总电荷为零，存在两个互不重叠的同味异号轻子对。\\
轻子对质量 & 所有同味异号对质量 $>5$；$50<m_{Z1}<106$，$12<m_{Z2}<115$。\\
质量窗与末态 & $105\leq m_{4\ell}<140$；随后要求两个电子和两个缪子。\\
\bottomrule
\end{tabularx}
\end{table}

\textbf{项目选择与文献的关系。} 沿用下载版既有核对：ATLAS 分析\cite{ATLAS13}支持接受度、
前三个动量尺度、配对、$Z$ 质量窗和横向碰撞参数的物理动机；本项目有以下改动：
\begin{itemize}
\item 使用统一 tight-ID；与文献宽松电子鉴别及详细缪子工作点的等价性未确认。
\item 动量允许取等号，第四个轻子统一 $p_T\geq7$；文献对电子、缪子的要求不同。
\item 两种隔离均要求小于 30\%；文献使用更细的隔离阈值、锥及修正规定。
\item 只保留恰好四个合格轻子的 $2e2\mu$ 道，未复现文献完整多候选与角分离程序。
\end{itemize}
教育示例\cite{HZZCuts,HZZAnalysis}支持部分计数、电荷、味道和动量规则，其隔离调用已注释。
Run-1 文献\cite{ATLASRun1}支持触发匹配原则。tight-ID 工作点和 $|z_0\sin\theta|<0.5$ mm 的确切外部依据仍待核实，
本报告把它们作为项目固定选择，不声称复现已发表测量。

\textbf{样本范围。} 信号 DSID 为 345060（\texttt{ggH125\_ZZ4lep}），背景为 363490（\texttt{llll}）。
背景的具体初态及 $Z/\gamma^*$ 贡献尚未核实，不能简称为纯夸克初态 $ZZ$。
发布前的精确预筛选和两份样本的生成器、簇射、探测器模拟配置仍待相应生产记录核实。
公开样本已经预筛选，本文无法据此计算相对全部生成事件的接受效率。

\clearpage
\reportsection{附录 B：归一化、来源与修订说明}
\textbf{物理归一化常数。} 沿用下载版及其固定版本参数表\cite{SampleInfo}：
\[
\begin{array}{c|cc}
 & \sigma_{\rm eff}\ (\mathrm{pb}) & S_{\rm gen}\\\hline
\text{信号} & 0.0060239 & 27881776.6536\\
\text{背景} & 1.2578 & 7538705.8077
\end{array}
\]
积分亮度 $\mathcal L=10^4\,\mathrm{pb}^{-1}=10\,\mathrm{fb}^{-1}$。
项目不另乘 $k$ 因子或过滤效率（配置均为 1），也未乘堆积、轻子及触发尺度因子。
有效截面中分支比与过滤效率的具体分解仍待核实，因此不能据此复现官方分析的产额。

\textbf{本次修订依据。} 基础文件为 Downloads 中的 \texttt{project-report-cn-v4.pdf}（10 页），
对应的仓库文本参照提交 \texttt{bb1d4cd}。本版沿用其中 ABCD、seed 42 的输入、训练说明、阈值、
事件组数、产额、区域占比，以及两份已跟踪图中的 ABCD 面板。原文件和其他报告均保留。
正文改为五步解释，删去输入组合比较与区间推断内容；教学阈值例子继续明确标注为假设数据。

\textbf{数值来自哪次运行？} 既有证据索引
\nolinkurl{paper/draft-v2/DRAFT001-EVIDENCE.md} 记录报告
\nolinkurl{runs/h4l-draft-v2-draft001/report/} 的身份为
\nolinkurl{9aab6d274a3e2d00d892c2d8fbf89b1e46f5b7a95d0bf0c085af90672ea36898}。
当前 Mac 中同名运行只记录到冻结，其冻结身份与该历史报告不同；本版不将二者视为同一份已完成结果。
本次核对了下载版文本、仓库协议和已跟踪图的哈希；本机没有这份历史完整报告，未重新审计其运行产物。
没有打开事件载荷、重新训练、评分、计算阈值或拟合。

\textbf{使用和解释范围。} 本研究属于教育与技术 MC 研究，探索性、非独立，不能作为 ATLAS/CMS 官方结果、
希格斯发现或物理测量。文献只提供方法与来源背景，不提供本报告的数值结果。
外部来源说明沿用下载版的核对记录，本次没有新增外部来源验证。CERN 发布说明将该集合限定为教育用途。

\begingroup\footnotesize
\renewcommand{\refname}{参考来源}
\begin{thebibliography}{9}
\bibitem{OpenData} ATLAS Collaboration, \emph{ATLAS 13 TeV samples collection at least four leptons (electron or muon), for 2020 Open Data release}.
\href{https://opendata.cern.ch/record/15005}{CERN Open Data record 15005};
\href{https://doi.org/10.7483/OPENDATA.ATLAS.2Y1T.TLGL}{doi:10.7483/OPENDATA.ATLAS.2Y1T.TLGL}.
\bibitem{ATLAS13} ATLAS Collaboration, \emph{Measurement of the Higgs boson coupling properties in the $H\to ZZ^*\to4\ell$ decay channel at $\sqrt{s}=13$ TeV with the ATLAS detector}, JHEP \textbf{03} (2018) 095.
\href{https://arxiv.org/pdf/1712.02304v2#page=8}{arXiv:1712.02304v2，第 5.1--5.2 节}。
\bibitem{HZZCuts} ATLAS Open Data outreach tools, \emph{HZZCuts.py}, revision \texttt{8ad2015be6d350060c3a183aff5802570bc9ada4}.
\href{https://github.com/atlas-outreach-data-tools/atlas-outreach-Python-uproot-framework-13tev/blob/8ad2015be6d350060c3a183aff5802570bc9ada4/HZZCuts.py#L2-L33}{固定版本源码}。
\bibitem{HZZAnalysis} ATLAS Open Data outreach tools, \emph{HZZAnalysis.py}, same revision as [3].
\href{https://github.com/atlas-outreach-data-tools/atlas-outreach-Python-uproot-framework-13tev/blob/8ad2015be6d350060c3a183aff5802570bc9ada4/HZZAnalysis.py#L140-L168}{固定版本源码}。
\bibitem{ATLASRun1} ATLAS Collaboration, \emph{Measurements of Higgs boson production and couplings in the four-lepton channel in pp collisions at center-of-mass energies of 7 and 8 TeV}, Phys. Rev. D \textbf{91}, 012006 (2015).
\href{https://arxiv.org/pdf/1408.5191v3#page=7}{arXiv:1408.5191v3，第 5.1 节}。
\bibitem{SampleInfo} ATLAS Open Data outreach tools, \emph{infofile.py}, same revision as [3].
\href{https://github.com/atlas-outreach-data-tools/atlas-outreach-Python-uproot-framework-13tev/blob/8ad2015be6d350060c3a183aff5802570bc9ada4/infofile.py}{固定版本样本参数表}。
\end{thebibliography}
\endgroup
\end{document}
